\scanpage{74}\mspage{I}
Dans le cas où $k$ est de car. $0$, on
peut supposer que $k$ est \deleted{de type fini
sur \unsure{$\mathbb Q$}}, le corps \unsure{des|de} fonctions d'un
schéma intègre $S$ de type fini sur $\mathbb Z$,
et que $f:X\unsure{-|\to}\unsure{Y|X}$ provient d'un
morphisme de $S$-préschémas $\bar f:\overline X\to\overline Y$,
intègres, $\overline Y$ simple sur $S$, et $\overline X$ \deleted{simple}
normal sur $S$, les fibres de l'un et
l'autre étant univ.$^{\mathrm t}$ connexes. \deleted{Alors pour
\unsure{tt|$\forall$} $y\in Y$, \unsure{s.c.|si} \unsure{$\mathrm{Ker}$} $H^1(\unsure{Y|X},\underline{O}_Y)\to H^1(X,\underline{O}_X)$
est de dim $r$, il existe un} on
\deleted{peut supposer qu'il existe un} On peut
supposer aussi $X$ et $Y$ coh. plats sur
$S$ en dim $1$, \deleted{enfin que} alors
\deleted{\unsure{cette}} \unsure{Ker|l'hom} $R^1\unsure{q_*|p_*}(\underline{O}_{\overline Y})\unsure{-|\to} R^1p_*(\underline{O}_{\overline X})$ est injectif,
\deleted{et que} (car il l'est en tt pt
par \unsure{ci-dessus dit|ci dessus dit}) alors \unsure{restreignant} $S$, si
le \unsure{conoyau|corang} de $H^1(Y,\underline{O}_Y)\unsure{-|\to}H^1(X,\underline{O}_X)$
est de dim $r$, la $=$ chose sera
vraie pour les fibres en \uline{tous} les \unsure{pts} \deleted{\unsure{de}}
$s\in S$, en particulier en \unsure{aux|des|au} pts de
% archive sheet number: 73
